\chapter{Conclusiones y trabajo futuro}\label{cap:conclusiones}

Este capítulo cierra la memoria repasando en qué medida se han alcanzado
los objetivos planteados en el capítulo~\ref{cap:introduccion}, sintetizando
los resultados más relevantes del proyecto, y proponiendo las líneas de
trabajo futuro, algunas de las cuales se han ido señalando a lo largo del documento.

\section{Cumplimiento de los objetivos}
\label{sec:conclusiones-objetivos}

Los once objetivos específicos planteados se han cumplido en su totalidad.

OE-01 a OE-05 el núcleo de la comparación de algoritmos: estudiar el
dataset, implementar ByteTrack y OC-SORT, ejecutarlos sobre la totalidad de
train y test, evaluarlos con las métricas estándar del campo y realizar un
estudio sistemático de hiperparámetros, se han cumplido de forma directa
y están documentados con resultados reales. OE-06 a OE-09
se han cumplido igualmente, con la salvedad, ya señalada en su
momento, de que OE-08 se resolvió como prueba de concepto sobre dos tramos
concretos de cámara estable y no como una solución general al movimiento
continuo de cámara del resto del dataset.
OE-10, documentar de forma
honesta las limitaciones y los enfoques descartados, no es un objetivo que
se cumpla en un capítulo concreto sino un criterio aplicado a todo el documento,
desde la tabla de problemas de entorno del
capítulo~\ref{cap:implementacion} hasta el propio capítulo~\ref{cap:uso-ia}.
OE-11, desarrollar una aplicación web como demostrador visual e
interactivo, se ha completado como un MVP funcional
fiel al diseño planificado en el
capítulo~\ref{cap:diseno}, aunque limitado a las tres secuencias de
demostración usadas a lo largo de la memoria y no al dataset completo.

\section{Resultados principales}
\label{sec:conclusiones-resultados}

Este proyecto como resultado ha dado una comparativa entre dos algoritmos de tracking: ByteTrack y OC-SORT. 
El resultado más relevante de este proyecto no es una cifra aislada, sino la comparativa entre ByteTrack y OC-SORT depende
de forma decisiva de qué variante de HOTA se use para medirla.
Bajo HOTA(0)  ByteTrack
toma ventaja a OC-SORT en el conjunto de test, sin em embargo, bajo HOTA integrada ese resultado se
invierte con claridad en ambos conjuntos, porque OC-SORT reutiliza las
detecciones del \emph{dataset} sin modificarlas mientras que ByteTrack las
interpola en su segunda ronda de asociación, lo que penaliza a ByteTrack en
los umbrales más estrictos que HOTA(0) no llega a evaluar. Son dos formas
igual de buenas para medir el mismo sistema. Que lleven a conclusiones
opuestas sobre qué algoritmo es mejor no es un detalle menor de este
proyecto, es uno de sus resultados.

Sobre la configuración óptima resultante, el pipeline de metadatos
ha demostrado ser aplicable sobre
secuencias reales de SoccerNet, con resultados numéricos consistentes,
por ejemplo, el porcentaje de posesión del equipo dominante en
SNMOT-116 se mantiene prácticamente igual pese a partir de
asociaciones de tracking distintas.

\section{Limitaciones}
\label{sec:conclusiones-limitaciones}

Las limitaciones de este proyecto se han ido documentando a lo largo del desarrollo en
los capítulos del proyecto, pero se resumen aquí:

\begin{itemize}
  \item El estudio de hiperparámetros se realiza en HOTA(0), no en
        la métrica integrada usada en la comparativa principal
        (sección~\ref{sec:pruebas-limitaciones}).
  \item La comparativa de HOTA integrada usa los valores por defecto de
        cada algoritmo, no sus configuraciones óptimas de hiperparámetros
        (sección~\ref{sec:pruebas-limitaciones}).
  \item La clasificación de equipos por K-means con $k=2$ asigna
        incorrectamente al árbitro a uno de los dos equipos
        (sección~\ref{sec:fundamentos-kmeans}).
  \item La homografía real se ha limitado a dos tramos cortos de cámara
        estable, no al dataset completo
        (sección~\ref{sec:implementacion-homografia}).
  \item La posesión de balón es una aproximación por proximidad espacial,
        no una estadística táctica real
        (sección~\ref{subsec:diseno-posesion}).
  \item La detección del balón se basa en un filtro geométrico sobre las
        detecciones del \emph{dataset}, no en un detector entrenado
        específicamente para esa clase.
  \item La aplicación web cubre únicamente las tres secuencias de
        demostración usadas en esta memoria, no el dataset completo para los metadatos.
\end{itemize}

\section{Trabajo futuro}
\label{sec:conclusiones-futuro}

De las limitaciones anteriores, y de los enfoques explorados y descartados
a lo largo del proyecto, se
derivan directamente las líneas de continuación más naturales de este
trabajo:

\begin{itemize}
  \item \textbf{Recalcular el estudio de hiperparámetros con HOTA
        integrada}: repetir el barrido completo de ambos algoritmos
        con la métrica integrada en
        vez de HOTA(0), y evaluar si la configuración óptima resultante
        coincide con la ya encontrada.
  \item \textbf{Evaluar las configuraciones óptimas con HOTA integrada}:
        confirmar si la ventaja de OC-SORT sobre ByteTrack observada con
        los valores por defecto de cada algoritmo se mantiene, se amplía o
        se reduce al usar sus configuraciones óptimas de hiperparámetros.
  \item \textbf{Incorporar algoritmos con re-identificación visual
        (Re-ID)}, como DeepSORT,
        aunque previsiblemente con un margen de mejora menor que en otros
        dominios de MOT, dada la escasa discriminabilidad visual entre
        compañeros de un mismo equipo.
  \item \textbf{Clasificación de equipos con un tercer grupo para el
        árbitro} ($k=3$, combinado con una heurística sobre el tamaño de
        cada grupo), para eliminar el sesgo que introduce actualmente en
        los mapas de calor por equipo y en la posesión de balón.
  \item \textbf{Homografía generalizada y dinámica}, que recalcule la
        matriz de proyección frame a frame para tramos con movimiento de
        cámara continuo, en lugar de limitarse a tramos concretos con
        cámara estática.
  \item \textbf{Extender la aplicación web al resto del dataset}, generando
        el paquete de metadatos y vídeos anotados para más secuencias o
        para su totalidad en lugar de las tres actuales, y explorando un
        detector de balón entrenado que sustituya al filtro geométrico
        actual como base de una posesión más fiable.
  \item \textbf{Tracking sobre un vídeo subido por el usuario}: permitir que
        el propio usuario suba un vídeo de fútbol arbitrario a la
        aplicación web y ejecutar el pipeline de tracking sobre él, en
        lugar de limitarse a las secuencias de demostración ya procesadas.
        A diferencia del resto de este proyecto, que parte de detecciones
        ya calculadas por SoccerNet (sección~\ref{sec:intro-problema}), esta
        vía exigiría incorporar primero un detector de objetos propio, ya
        que un vídeo arbitrario no viene con \texttt{det.txt}. Es
        técnicamente viable con el enfoque ya desarrollado, pero queda
        como línea futura.
  \item \textbf{Participación real en el \emph{leaderboard} oficial de
        SoccerNet}: el dataset incluye un tercer split, \emph{challenge},
        sin \texttt{gt.txt},
        pensado precisamente para enviar resultados al servidor oficial de
        evaluación y obtener una comparación pública y directa con otros
        equipos, en la línea de los ya publicados en ediciones anteriores
        del reto \citep{cioppa2023soccernetchallenges}, ofreciendo esta vez
        una cifra oficial en lugar de la referencia de contexto de la
        sección~\ref{sec:fundamentos-relacionados}.
\end{itemize}

\section{Reflexión final}
\label{sec:conclusiones-reflexion}

Este proyecto partió de una pregunta relativamente acotada
 y ha terminado abarcando un sistema completo, desde el
seguimiento de jugadores hasta la extracción de metadatos deportivos
interpretables y su exposición en una aplicación web. La constatación de que la propia pregunta
\enquote{qué algoritmo es mejor}depende de una decisión metodológica, la métrica
elegida para responderla, que resultó menos trivial de lo que parecía al
principio del proyecto.
